The topic of this thesis is statistical inference for Richardson--Romberg
averaged constant-stepsize linear stochastic approximation with Markovian
noise. The object of study is the Polyak--Ruppert averaged LSA trajectory,
while the subject is the distribution of scalar projections of the RR
estimator and its covariance normalization. The methods used in the work are
deterministic error decompositions, Lyapunov contraction estimates,
Markov-chain Poisson equations, martingale Berry--Esseen inequalities,
Richardson--Romberg weight analysis, and finite-start burn-in transfer. The
problem is relevant because constant-stepsize stochastic approximation is
widely used in reinforcement learning and stochastic optimization, but its
steady-state bias complicates confidence interval construction under
dependent data. The goal is to obtain a non-asymptotic normal approximation
for the averaged RR statistic. The main tasks are to identify the leading
linearized term, control RR weights and misadjustment terms, transfer the
stationary theorem to deterministic starts, and check the practical behavior
of the resulting confidence intervals. The main result is a balanced-scale
Berry--Esseen bound with covariance target $\Sigma_\infty$. The thesis is
organized as follows: the first chapters develop the deterministic and
stationary proof ingredients, the burn-in chapter transfers the result to
finite starts, the experiment chapter records numerical diagnostics, and the
appendices collect key objects and external inputs.

Stochastic approximation (SA) algorithms are a cornerstone of modern
computational statistics, optimization, and reinforcement learning.
Introduced by Robbins and Monro (1951), these iterative procedures provide a
principled way to find roots of equations or optimize objectives when only
noisy observations are available. A particularly important subclass is the
\emph{linear stochastic approximation} (LSA) algorithm, which arises
naturally in temporal-difference (TD) learning, policy evaluation, and
stochastic gradient descent for linear models.

In this work, we study the LSA recursion with a \emph{constant step size}
$\alpha>0$:
\begin{equation}
  \theta_k^{(\alpha)}
  = \theta_{k-1}^{(\alpha)}
    - \alpha\bigl\{A(Z_k)\theta_{k-1}^{(\alpha)}-b(Z_k)\bigr\},
  \qquad k\geq 1.
  \label{eq:lsa}
\end{equation}
Here $\{Z_k\}_{k\in\mathbb{N}}$ is a time-homogeneous Markov chain on a
measurable space $(\mathsf{Z},\mathcal{Z})$ with transition kernel $\mathsf{Q}$
and unique invariant distribution $\pi$. The mappings
$A\colon\mathsf{Z}\to\mathbb{R}^{d\times d}$ and
$b\colon\mathsf{Z}\to\mathbb{R}^d$ are measurable functions satisfying
\[
  \overline{A}:=\int_{\mathsf{Z}}A(z)\,\mathrm{d}\pi(z),
  \qquad
  \overline{b}:=\int_{\mathsf{Z}}b(z)\,\mathrm{d}\pi(z).
\]
We use the stability convention that $-\overline{A}$ is Hurwitz, equivalently
all eigenvalues of $\overline{A}$ have strictly positive real parts. Then the
target parameter $\theta^*=\overline{A}^{-1}\overline{b}$ is uniquely defined.

The \emph{Polyak--Ruppert averaging} procedure (Polyak, 1990; Ruppert, 1988)
provides an effective variance reduction technique. Given a burn-in period
$n_0\geq 0$, the averaged iterate is defined as
\begin{equation}
  \overline{\theta}_n^{(\alpha)}
  = \frac{1}{n-n_0}\sum_{k=n_0}^{n-1}\theta_k^{(\alpha)}.
  \label{eq:pr-average}
\end{equation}
